\chapter{Code Repository Structure}

\section{Repository Overview}

The complete source code for this thesis is available at: \texttt{[GitHub URL]}

\section{Directory Structure}

\begin{verbatim}
edge-ai-vineyard-monitoring/
├── dd_cnn/                          # Disease classification training
│   ├── Model_training/
│   │   ├── MobileNetV2.ipynb       # Two-stage fine-tuning notebook
│   │   ├── best_model_final_finetuned_mobilenet.pth
│   │   ├── best_model_stage1_mobilenet.pth
│   │   └── data/                   # Training dataset
│   └── dataset/
│       └── grape-disease/          # Grape leaf disease images
│
├── YOLO/                            # Leaf detection training
│   ├── YOLO11_ESP32_TFLite_Export.ipynb
│   ├── yolo11n.pt                  # YOLOv11 nano model
│   ├── Dataset.v1.yolov11/         # YOLO format dataset
│   └── esp32_deployment_package/   # Deployment files
│
├── esp-detection/                   # ESP32-S3 firmware
│   └── esp-detection/
│       ├── main/
│       │   ├── main.cpp            # Main inference pipeline
│       │   ├── camera_handler.cpp  # OV3660 camera driver
│       │   ├── detection_model.cpp # ESPDet-Pico inference
│       │   ├── classification_model.cpp # MobileNetV2 inference
│       │   └── aggregation.cpp     # Multi-leaf aggregation
│       ├── models/                  # Quantized model files (.espdl)
│       └── CMakeLists.txt
│
└── Documentation/                   # Thesis LaTeX files
    ├── main.tex
    ├── chapters/
    ├── appendices/
    └── Bibliography.bib
\end{verbatim}

\section{Key Files Description}

\subsection{Training Files}
\begin{itemize}
    \item \texttt{MobileNetV2.ipynb}: Two-stage fine-tuning pipeline with diagnostic cells
    \item \texttt{YOLO11\_ESP32\_TFLite\_Export.ipynb}: YOLO model training and quantization
    \item \texttt{yolo11n.pt}: Pre-trained YOLOv11 nano weights
\end{itemize}

\subsection{Firmware Files}
\begin{itemize}
    \item \texttt{main.cpp}: End-to-end inference pipeline orchestration
    \item \texttt{camera\_handler.cpp}: DMA buffer management and JPEG capture
    \item \texttt{detection\_model.cpp}: ESPDet-Pico inference with 416×320 input
    \item \texttt{classification\_model.cpp}: MobileNetV2 inference with 128×128 input
    \item \texttt{aggregation.cpp}: Hybrid weighted aggregation algorithm
\end{itemize}

\section{Building and Running}

\subsection{Training Models}
\begin{lstlisting}[language=bash]
# Install dependencies
pip install torch torchvision ultralytics

# Train MobileNetV2 classification model
jupyter notebook dd_cnn/Model_training/MobileNetV2.ipynb

# Train YOLO detection model
jupyter notebook YOLO/YOLO11_ESP32_TFLite_Export.ipynb
\end{lstlisting}

\subsection{Flashing ESP32-S3}
\begin{lstlisting}[language=bash]
# Setup ESP-IDF environment
cd esp-detection/esp-detection
idf.py set-target esp32s3

# Build and flash
idf.py build
idf.py -p /dev/ttyUSB0 flash monitor
\end{lstlisting}

\section{Repository Access}

For detailed build instructions, see Appendix C.
